\begin{abstract}
Reasoning models check, revise, decompose, and explore alternatives. How are these shared operations organized for different tasks, and what happens when that organization is changed? Thought Atlas analyzes 21,803 reasoning traces from seven models across eight benchmark domains. Operation prevalence improves held-out domain prediction beyond length in every tested model and annotation group; temporal allocation adds information beyond prevalence. Mathematical traces increasingly emphasize verification, whereas idea-generation traces emphasize alternatives and candidate revision. Matched instruction experiments show that changing task demand can alter these profiles beyond repeated-generation noise, with substantial model heterogeneity. On Qwen3.5-9B, activation interventions constructed from domain profiles produce asymmetric transfer: a math-associated profile improves science and planning accuracy at a restricted token budget, while reverse transfers can be harmful. These benefits largely diminish at a longer cutoff. A reward-trained residual bias produces a behavioral shift similar to an atlas-derived intervention, yet differs in early forced-answer performance. Objective-task benefits can also accompany lower judged originality on an open task. The findings connect the task-dependent organization of a shared behavioral repertoire to the benefits and limits of reasoning control: a profile's usefulness depends on the task, budget, and objective.
\end{abstract}
